\documentclass[12pt]{ctexart}
\usepackage{a4}
\usepackage{graphicx}
\usepackage{geometry}
\usepackage{amsmath}
\usepackage{mathtools}
\usepackage{amssymb}
\geometry{a4paper,left=2cm,right=2cm,top=2.5cm,bottom=2.5cm}

\title{几何基础作业6}
\author{无名}
\date{\today}

\begin{document}
\maketitle

\section{证明“有理数之树”包含所有的正有理数，且每个有理数都只出现一次}

\begin{figure}[htb]
    \begin{center}
        \includegraphics[width=0.8\linewidth]{1.png}
        \caption{}
    \end{center}
\end{figure}
\subsection{定义}
定义：对于节点$\frac{a}{b}$，其左侧子节点为$\frac{a+b}{b}$，右侧为$\frac{a}{a+b}$
\subsection{证明树内所有分数都最简}
考虑所有的正有理数为$\frac{m}{n}\hspace{2em}m,n\in \mathbb{N} ^+$
那么显然可见，作为最简形式，$m,n$必须互为质数，即需要证明树中会出现所有最简分数。

考虑$\frac{a}{b}$最简，那么$a,b$互质。

若$\frac{a+b}{b}$非最简，那么假设$d=gcd(a,a+b)$，从而$d \vert a $且$d \vert a+b$，故一定会有$d\vert (a+b)-a=b$

这与$a,b$互质矛盾，显然不成立。

同理，右侧$\frac{a}{a+b}$显然也该最简。

而有理数之树的第一层为$\frac{1}{1}$，那么由$\frac{1}{1}$最简可以得到树内所有分数都最简
\subsection{证明最简分数一定能被找到且唯一}
假设$\frac{m}{n}$在树中，那么

若$m>n$那么它只能是上一级的左侧子节点，其母节点为$\frac{m-n}{n}$

若$m<n$那么它只能是上一级的右侧子节点，母节点为$\frac{m}{n-m}$

若$m=n$那么它只能是$\frac{1}{1}$

故所有的树中有理数都有唯一的路径到$\frac{1}{1}$，那么只能出现一次

由于起始节点为$\frac{1}{1}$，则m,n任一个取1可以使得另外一个不断累加从而最终取值属于$\mathbb{N^+}$，根据定义这就能构造出任意的有理数$\frac{m}{n}$

证毕
\end{document}